\documentclass{article}
\usepackage{amsmath,amssymb,amsthm}
\usepackage{algorithm2e}
\usepackage{enumitem}
\usepackage{hyperref}
\usepackage{geometry}
\geometry{margin=1in}
\newtheorem{theorem}{Theorem}
\newtheorem{lemma}{Lemma}
\newcommand{\prox}{\operatorname{prox}}
\newcommand{\R}{\mathbb{R}}
\newcommand{\E}{\mathbb{E}}

\begin{document}

\title{Convergence Analysis of the Proximal Gradient Method}
\author{}
\date{}
\maketitle

\tableofcontents

\newpage

%%%%%%%%%%%%%%%%%%%%%%%%%%%%%%%%%%%%%%%%%%%%%%%%%%%%%%%%%%%%%%%%%
\section*{Algorithm Name}
\addcontentsline{toc}{section}{Algorithm Name}
\begin{center}
\textbf{Proximal Gradient Method (Forward–Backward Splitting)}
\end{center}
%%%%%%%%%%%%%%%%%%%%%%%%%%%%%%%%%%%%%%%%%%%%%%%%%%%%%%%%%%%%%%%%%

\section*{Mathematical Formulation}
\addcontentsline{toc}{section}{Mathematical Formulation}
Let 
\[
F(x):=f(x)+g(x),\qquad x\in\R^{d},
\]
where 
\begin{itemize}[noitemsep]
\item $f:\R^{d}\to\R$ is convex and differentiable with $L$‑Lipschitz continuous gradient,
\item $g:\R^{d}\to\R\cup\{+\infty\}$ is proper, closed, convex (possibly non‑smooth).
\end{itemize}
Given a stepsize $\eta>0$, the proximal gradient iteration is
\begin{equation}
\label{eq:update}
x_{k+1}= \prox_{\eta g}\!\bigl(x_{k}-\eta\nabla f(x_{k})\bigr),
\qquad k=0,1,2,\dots
\end{equation}
where the proximal operator of $\eta g$ is defined by
\[
\prox_{\eta g}(v):=\arg\min_{u\in\R^{d}}\Bigl\{ g(u)+\frac{1}{2\eta}\|u-v\|^{2}\Bigr\}.
\]

%%%%%%%%%%%%%%%%%%%%%%%%%%%%%%%%%%%%%%%%%%%%%%%%%%%%%%%%%%%%%%%%%
\section*{Pseudocode}
\addcontentsline{toc}{section}{Pseudocode}
\begin{algorithm}[H]
\SetAlgoLined
\KwIn{Initial point $x_{0}\in\R^{d}$, stepsize $\eta\in(0,2/L)$, maximal iterations $K$}
\For{$k\leftarrow 0$ \KwTo $K-1$}{
    $v_{k}\gets x_{k}-\eta\nabla f(x_{k})$\;
    $x_{k+1}\gets \prox_{\eta g}(v_{k})$\;
}
\KwOut{Sequence $\{x_{k}\}_{k=0}^{K}$}
\caption{Proximal Gradient Method}
\end{algorithm}
%%%%%%%%%%%%%%%%%%%%%%%%%%%%%%%%%%%%%%%%%%%%%%%%%%%%%%%%%%%%%%%%%

\section*{Dry Run Example}
\addcontentsline{toc}{section}{Dry Run Example}
Consider the scalar problem
\[
f(w)=(w-3)^{2},\qquad g\equiv0,
\]
so that $\prox_{\eta g}(v)=v$.  Take $w_{0}=0$, stepsize $\eta=0.1$, and run three iterations.

\begin{center}
\begin{tabular}{c|c|c|c}
Iteration $k$ & $\nabla f(w_{k})$ & $v_{k}=w_{k}-\eta\nabla f(w_{k})$ & $w_{k+1}=v_{k}$ \\ \hline
0 & $2(w_{0}-3)=-6$ & $0-0.1(-6)=0.6$ & $0.6$ \\ 
1 & $2(0.6-3)=-4.8$ & $0.6-0.1(-4.8)=1.08$ & $1.08$ \\ 
2 & $2(1.08-3)=-3.84$ & $1.08-0.1(-3.84)=1.464$ & $1.464$ \\ 
\end{tabular}
\end{center}
The corresponding objective values $F(w_{k})=f(w_{k})$ are
\[
f(0)=9,\quad f(0.6)=5.76,\quad f(1.08)=3.6864,\quad f(1.464)=2.3415.
\]
We see a monotone decrease toward the optimum $w^{\star}=3$.

%%%%%%%%%%%%%%%%%%%%%%%%%%%%%%%%%%%%%%%%%%%%%%%%%%%%%%%%%%%%%%%%%
\section*{Assumptions}
\addcontentsline{toc}{section}{Assumptions}
\begin{enumerate}[label={A\arabic*}:, leftmargin=*, noitemsep]
\item \textbf{Smoothness of $f$.} $\nabla f$ is $L$‑Lipschitz:
\[
\|\nabla f(x)-\nabla f(y)\|\le L\|x-y\|,\qquad\forall x,y\in\R^{d}.
\tag{A1}
\]
\item \textbf{Convexity of $f$ and $g$.}
\[
f(\theta x+(1-\theta)y)\le \theta f(x)+(1-\theta)f(y),\quad
g(\theta x+(1-\theta)y)\le \theta g(x)+(1-\theta)g(y)
\]
for all $x,y$ and $\theta\in[0,1]$.
\tag{A2}
\item \textbf{Existence of a minimizer.} The set
\[
\mathcal{X}^{\star}:=\arg\min_{x\in\R^{d}}F(x)
\]
is non‑empty, and we denote $F^{\star}=F(x^{\star})$ for any $x^{\star}\in\mathcal{X}^{\star}$.
\tag{A3}
\item \textbf{Stepsize.} The constant stepsize satisfies
\[
0<\eta<\frac{2}{L}.
\tag{A4}
\]
\end{enumerate}

%%%%%%%%%%%%%%%%%%%%%%%%%%%%%%%%%%%%%%%%%%%%%%%%%%%%%%%%%%%%%%%%%
\section*{Main Convergence Theorem}
\addcontentsline{toc}{section}{Main Convergence Theorem}
\begin{theorem}[Suboptimality $O(1/k)$]
\label{thm:rate}
Under Assumptions A1–A4, the iterates generated by \eqref{eq:update} satisfy for all $k\ge 1$
\[
F(x_{k})-F^{\star}\;\le\;\frac{\|x_{0}-x^{\star}\|^{2}}{2\eta\,k},
\qquad\forall\,x^{\star}\in\mathcal{X}^{\star}.
\]
Consequently $F(x_{k})\to F^{\star}$ and any limit point of $\{x_{k}\}$ belongs to $\mathcal{X}^{\star}$.
\end{theorem}
The proof relies solely on the non‑expansiveness of the proximal operator and the descent lemma for $f$.

%%%%%%%%%%%%%%%%%%%%%%%%%%%%%%%%%%%%%%%%%%%%%%%%%%%%%%%%%%%%%%%%%
\section*{Theoretical Properties}
\addcontentsline{toc}{section}{Theoretical Properties}
\begin{enumerate}[label={P\arabic*}:, leftmargin=*, noitemsep]
\item \textbf{Stability.} Because $\prox_{\eta g}$ is firmly non‑expansive,
\[
\|\prox_{\eta g}(u)-\prox_{\eta g}(v)\|^{2}\le\langle u-v,\prox_{\eta g}(u)-\prox_{\eta g}(v)\rangle,
\]
the algorithm is Lyapunov stable for any $\eta\in(0,2/L)$.
\item \textbf{Hyperparameter Sensitivity.}
The bound in Theorem~\ref{thm:rate} is monotone decreasing in $\eta$ up to the maximal admissible value $2/L$, indicating that larger steps (still below $2/L$) accelerate convergence.
\item \textbf{Comparison to Gradient Descent.}
When $g\equiv0$, $\prox_{\eta g}$ reduces to the identity and the bound becomes the classical $O(1/k)$ rate for gradient descent on smooth convex functions.
\item \textbf{Extension to Composite Strong Convexity.}
If $F$ is $\mu$‑strongly convex, the same proof yields a linear rate
\[
\|x_{k}-x^{\star}\|^{2}\le\bigl(1-\eta\mu\bigr)^{k}\|x_{0}-x^{\star}\|^{2}.
\]
\end{enumerate}

%%%%%%%%%%%%%%%%%%%%%%%%%%%%%%%%%%%%%%%%%%%%%%%%%%%%%%%%%%%%%%%%%
\section*{Preliminaries (Definitions and Lemmas)}
\addcontentsline{toc}{section}{Preliminaries (Definitions and Lemmas)}
\begin{lemma}[Descent Lemma for $f$]
\label{lem:descent}
If $\nabla f$ is $L$‑Lipschitz, then for any $x,y\in\R^{d}$
\[
f(y)\le f(x)+\langle\nabla f(x),y-x\rangle+\frac{L}{2}\|y-x\|^{2}.
\tag{L1}
\]
\end{lemma}
\begin{proof}
Standard consequence of the integral form of the gradient; omitted for brevity. \qedhere
\end{proof}

\begin{lemma}[Firm Non‑Expansiveness of $\prox_{\eta g}$]
\label{lem:firm}
For any $u,v\in\R^{d}$,
\[
\|\prox_{\eta g}(u)-\prox_{\eta g}(v)\|^{2}
\le\langle u-v,\prox_{\eta g}(u)-\prox_{\eta g}(v)\rangle.
\tag{L2}
\]
\end{lemma}
\begin{proof}
Follows from the definition of the proximal operator as the resolvent of the sub‑differential $\partial g$, which is maximally monotone. \qedhere
\end{proof}

%%%%%%%%%%%%%%%%%%%%%%%%%%%%%%%%%%%%%%%%%%%%%%%%%%%%%%%%%%%%%%%%%
\section*{Main Proof (Step‑by‑Step Derivation)}
\addcontentsline{toc}{section}{Main Proof (Step-by-Step Derivation)}
We prove Theorem~\ref{thm:rate}.  Throughout the proof $x^{\star}\in\mathcal{X}^{\star}$ is fixed.

\noindent\textbf{Step 1.} Write the optimality condition of the proximal step.  
From the definition of $\prox_{\eta g}$,
\[
x_{k+1}= \prox_{\eta g}(v_{k}),\qquad v_{k}=x_{k}-\eta\nabla f(x_{k}),
\]
we have (see Lemma~\ref{lem:firm})
\[
\frac{1}{\eta}\bigl(v_{k}-x_{k+1}\bigr)\in\partial g(x_{k+1}).
\tag{Step 1}
\]

\noindent\textbf{Step 2.} Use convexity of $g$ with the sub‑gradient from Step 1.
For any $x^{\star}$,
\[
g(x^{\star})\ge g(x_{k+1})+\left\langle\frac{1}{\eta}(v_{k}-x_{k+1}),\,x^{\star}-x_{k+1}\right\rangle.
\tag{Step 2}
\]

\noindent\textbf{Step 3.} Rearrange Step 2 and substitute $v_{k}=x_{k}-\eta\nabla f(x_{k})$:
\[
\begin{aligned}
g(x_{k+1})-g(x^{\star})
&\le\frac{1}{\eta}\langle x_{k+1}-v_{k},\,x_{k+1}-x^{\star}\rangle\\
&=\frac{1}{\eta}\langle x_{k+1}-x_{k}+\eta\nabla f(x_{k}),\,x_{k+1}-x^{\star}\rangle\\
&=\frac{1}{\eta}\langle x_{k+1}-x_{k},\,x_{k+1}-x^{\star}\rangle
   +\langle\nabla f(x_{k}),\,x_{k+1}-x^{\star}\rangle .
\end{aligned}
\tag{Step 3}
\]

\noindent\textbf{Step 4.} Expand the first inner product using the identity
$\langle a-b,c-b\rangle=\tfrac12(\|a-b\|^{2}+\|c-b\|^{2}-\|a-c\|^{2})$:
\[
\frac{1}{\eta}\langle x_{k+1}-x_{k},\,x_{k+1}-x^{\star}\rangle
=
\frac{1}{2\eta}\Bigl(\|x_{k}-x^{\star}\|^{2}
-\|x_{k+1}-x^{\star}\|^{2}
-\|x_{k+1}-x_{k}\|^{2}\Bigr).
\tag{Step 4}
\]

\noindent\textbf{Step 5.} Combine Step 3 and Step 4:
\[
\begin{aligned}
g(x_{k+1})-g(x^{\star})
&\le
\frac{1}{2\eta}\bigl(\|x_{k}-x^{\star}\|^{2}
-\|x_{k+1}-x^{\star}\|^{2}
-\|x_{k+1}-x_{k}\|^{2}\bigr)\\
&\quad +\langle\nabla f(x_{k}),\,x_{k+1}-x^{\star}\rangle .
\end{aligned}
\tag{Step 5}
\]

\noindent\textbf{Step 6.} Add $f(x_{k+1})-f(x^{\star})$ to both sides.
Using convexity of $f$,
\[
f(x_{k+1})-f(x^{\star})\le \langle\nabla f(x_{k}),\,x_{k+1}-x^{\star}\rangle
      +\frac{L}{2}\|x_{k+1}-x_{k}\|^{2},
\tag{Step 6}
\]
where the last inequality follows from Lemma~\ref{lem:descent} (L1) with $y=x_{k+1}$.

\noindent\textbf{Step 7.} Sum the inequalities from Step 5 and Step 6 to obtain a bound on the composite objective $F=f+g$:
\[
\begin{aligned}
F(x_{k+1})-F^{\star}
&= \bigl[f(x_{k+1})-f(x^{\star})\bigr]+\bigl[g(x_{k+1})-g(x^{\star})\bigr]\\
&\le
\frac{1}{2\eta}\bigl(\|x_{k}-x^{\star}\|^{2}
-\|x_{k+1}-x^{\star}\|^{2}
-\|x_{k+1}-x_{k}\|^{2}\bigr)\\
&\quad +\frac{L}{2}\|x_{k+1}-x_{k}\|^{2}.
\end{aligned}
\tag{Step 7}
\]

\noindent\textbf{Step 8.} Collect the terms containing $\|x_{k+1}-x_{k}\|^{2}$:
\[
-\frac{1}{2\eta}\|x_{k+1}-x_{k}\|^{2}
+\frac{L}{2}\|x_{k+1}-x_{k}\|^{2}
= -\frac{1}{2}\Bigl(\frac{1}{\eta}-L\Bigr)\|x_{k+1}-x_{k}\|^{2}.
\tag{Step 8}
\]

\noindent\textbf{Step 9.} Under Assumption A4, $\eta<2/L$, thus $\frac{1}{\eta}>L/2$ and in particular
\[
\frac{1}{\eta}-L>0.
\]
Consequently the term in Step 8 is non‑positive, and we can drop it to get a simpler inequality:
\[
F(x_{k+1})-F^{\star}
\le
\frac{1}{2\eta}\bigl(\|x_{k}-x^{\star}\|^{2}
-\|x_{k+1}-x^{\star}\|^{2}\bigr).
\tag{Step 9}
\]

\noindent\textbf{Step 10.} Rearrange Step 9:
\[
2\eta\bigl(F(x_{k+1})-F^{\star}\bigr)
\le
\|x_{k}-x^{\star}\|^{2}
-\|x_{k+1}-x^{\star}\|^{2}.
\tag{Step 10}
\]

\noindent\textbf{Step 11.} Sum the inequality of Step 10 for $k=0,\dots, K-1$ (telescoping sum):
\[
2\eta\sum_{k=0}^{K-1}\bigl(F(x_{k+1})-F^{\star}\bigr)
\le
\|x_{0}-x^{\star}\|^{2}
-\|x_{K}-x^{\star}\|^{2}
\le
\|x_{0}-x^{\star}\|^{2}.
\tag{Step 11}
\]

\noindent\textbf{Step 12.} Since each term $F(x_{k+1})-F^{\star}\ge0$, the average is bounded:
\[
\frac{1}{K}\sum_{k=0}^{K-1}\bigl(F(x_{k+1})-F^{\star}\bigr)
\le
\frac{\|x_{0}-x^{\star}\|^{2}}{2\eta K}.
\tag{Step 12}
\]

\noindent\textbf{Step 13.} Because $F(x_{k})$ is non‑increasing (Step 9), we have for any $k\ge1$
\[
F(x_{k})-F^{\star}\le
\frac{1}{k}\sum_{j=0}^{k-1}\bigl(F(x_{j+1})-F^{\star}\bigr)
\le\frac{\|x_{0}-x^{\star}\|^{2}}{2\eta k}.
\tag{Step 13}
\]

\noindent\textbf{Conclusion.} Step 13 is exactly the bound claimed in Theorem~\ref{thm:rate}.  Moreover,
$F(x_{k})\to F^{\star}$, and the boundedness of $\{x_{k}\}$ together with the optimality condition of the proximal step implies any cluster point belongs to $\mathcal{X}^{\star}$. \qed

%%%%%%%%%%%%%%%%%%%%%%%%%%%%%%%%%%%%%%%%%%%%%%%%%%%%%%%%%%%%%%%%%
\section*{Rate Derivation (With Constants)}
\addcontentsline{toc}{section}{Rate Derivation (With Constants)}
The constant appearing in the $O(1/k)$ rate is $\frac{\|x_{0}-x^{\star}\|^{2}}{2\eta}$.  Choosing the largest admissible stepsize $\eta=1/L$ (which satisfies A4) yields the tightest bound:
\[
F(x_{k})-F^{\star}\le \frac{L\|x_{0}-x^{\star}\|^{2}}{2k}.
\]

%%%%%%%%%%%%%%%%%%%%%%%%%%%%%%%%%%%%%%%%%%%%%%%%%%%%%%%%%%%%%%%%%
\section*{Special Cases}
\addcontentsline{toc}{section}{Special Cases}
\begin{itemize}[noitemsep]
\item \textbf{Pure Gradient Descent.} If $g\equiv0$, $\prox_{\eta g}=\mathrm{Id}$ and the analysis reduces to the classical gradient descent proof.
\item \textbf{Indicator Function.} Let $g=\iota_{\mathcal{C}}$ be the indicator of a closed convex set $\mathcal{C}$. Then $\prox_{\eta g}$ is the Euclidean projection onto $\mathcal{C}$, and the method becomes projected gradient descent with the same $O(1/k)$ guarantee.
\item \textbf{Strong Convexity.} If $F$ is $\mu$‑strongly convex, a simple modification of Step 9 (using the inequality $\langle \nabla f(x_{k})- \nabla f(x^{\star}), x_{k}-x^{\star}\rangle\ge\mu\|x_{k}-x^{\star}\|^{2}$) yields a linear convergence rate $ \|x_{k}-x^{\star}\|^{2}\le(1-\eta\mu)^{k}\|x_{0}-x^{\star}\|^{2}$.
\end{itemize}

%%%%%%%%%%%%%%%%%%%%%%%%%%%%%%%%%%%%%%%%%%%%%%%%%%%%%%%%%%%%%%%%%
\section*{Discussion}
\addcontentsline{toc}{section}{Discussion}
The proof above showcases how the non‑expansiveness (firm non‑expansiveness) of the proximal operator replaces the role of the Euclidean projection in projected gradient methods.  The only requirement on the stepsize is the standard $0<\eta<2/L$, identical to gradient descent, which highlights that the presence of a non‑smooth term $g$ does not worsen the convergence speed under convexity.  Extensions to stochastic gradients, accelerated variants, or variable stepsizes follow the same template with additional martingale concentration arguments.

\end{document}